\documentclass[11pt,a4paper]{article}
\usepackage[utf8]{inputenc}
\usepackage[italian]{babel}
\usepackage{geometry}
\usepackage{graphicx}
\usepackage{hyperref}
\usepackage{listings}
\usepackage{xcolor}
\usepackage{amsmath}
\usepackage{enumitem}
\usepackage{tikz}
\usepackage{booktabs}
\usepackage{longtable}
\usepackage{array}
\usepackage{multirow}

\geometry{margin=2.5cm}

% Configurazione codice
\lstset{
    language=JavaScript,
    basicstyle=\ttfamily\small,
    keywordstyle=\color{blue}\bfseries,
    stringstyle=\color{red},
    commentstyle=\color{green!60!black},
    numberstyle=\tiny\color{gray},
    numbers=left,
    numbersep=5pt,
    frame=single,
    breaklines=true,
    showstringspaces=false,
    tabsize=2
}

\title{\textbf{Architettura del Sistema ExtraTracker}\\
\large Documentazione Tecnica Completa}
\author{Documentazione Sistema}
\date{\today}

\begin{document}

\maketitle
\tableofcontents
\newpage

\section{Introduzione}

Questo documento descrive in modo esaustivo l'architettura del sistema \textbf{ExtraTracker}, un'applicazione Full-Stack per la gestione delle ore di lavoro, il tracciamento degli straordinari, la gestione degli obiettivi personali e lo studio con flashcard.

Il sistema è progettato seguendo i principi di:
\begin{itemize}
    \item \textbf{Multi-tenancy}: Isolamento completo dei dati per utente
    \item \textbf{Sicurezza}: Autenticazione JWT, rate limiting, protezione CSRF/XSS
    \item \textbf{Scalabilità}: Architettura modulare, caching con Redis, ottimizzazioni database
    \item \textbf{Manutenibilità}: Separazione delle responsabilità, pattern MVC, codice tipizzato
\end{itemize}

\section{Architettura Generale}

\subsection{Visione d'Insieme}

Il sistema segue un'architettura \textbf{Client-Server} disaccoppiata con tre livelli principali:

\begin{enumerate}
    \item \textbf{Frontend (Client)}: React 19 + TypeScript, Single Page Application (SPA)
    \item \textbf{Backend (Server)}: Node.js + Express.js, REST API
    \item \textbf{Database}: MongoDB (dati persistenti) + Redis (cache e sessioni)
\end{enumerate}

\subsection{Stack Tecnologico}

\subsubsection{Frontend}
\begin{itemize}
    \item \textbf{Framework}: React 19 con Vite
    \item \textbf{Linguaggio}: TypeScript
    \item \textbf{Styling}: TailwindCSS
    \item \textbf{State Management}: React Context API
    \item \textbf{Routing}: React Router v7
    \item \textbf{Animazioni}: Framer Motion
    \item \textbf{HTTP Client}: Axios con interceptors personalizzati
\end{itemize}

\subsubsection{Backend}
\begin{itemize}
    \item \textbf{Runtime}: Node.js
    \item \textbf{Framework}: Express.js
    \item \textbf{ODM}: Mongoose (MongoDB)
    \item \textbf{Autenticazione}: JWT (Access Token + Refresh Token)
    \item \textbf{Validazione}: Zod
    \item \textbf{Sicurezza}: Helmet, CORS, Rate Limiting, HPP
    \item \textbf{Password Hashing}: Argon2id
\end{itemize}

\subsubsection{Database e Servizi}
\begin{itemize}
    \item \textbf{Database Primario}: MongoDB (dati utente, progetti, obiettivi, worklog)
    \item \textbf{Cache/Sessioni}: Redis (rate limiting distribuito, blacklist token)
    \item \textbf{Orchestrazione}: Docker Compose
\end{itemize}

\section{Componenti Frontend}

Il frontend è organizzato seguendo un pattern \textbf{Feature-First}, dove ogni feature è un modulo autonomo con i propri componenti, servizi, context e tipi.

\subsection{Struttura delle Features}

\subsubsection{1. Autenticazione (\texttt{features/auth})}

\textbf{Problema risolto}: Gestione completa del ciclo di vita dell'autenticazione utente.

\textbf{Componenti}:
\begin{itemize}
    \item \texttt{AuthContext.tsx}: Context provider per stato autenticazione globale
    \item \texttt{LoginPage.tsx}: Pagina di login
    \item \texttt{RegisterPage.tsx}: Pagina di registrazione
    \item \texttt{ForgotPasswordPage.tsx}: Recupero password
    \item \texttt{ResetPasswordPage.tsx}: Reset password con token
    \item \texttt{VerifyEmailPage.tsx}: Verifica email
\end{itemize}

\textbf{Servizi}:
\begin{itemize}
    \item \texttt{authService.ts}: Wrapper per chiamate API di autenticazione
\end{itemize}

\textbf{Comunicazione}:
\begin{itemize}
    \item \textbf{API utilizzate}: \texttt{POST /api/auth/login}, \texttt{POST /api/auth/register}, \texttt{POST /api/auth/refresh}, \texttt{GET /api/auth/check}
    \item \textbf{Dati scambiati}: Credenziali (email/password), JWT tokens (in cookie HttpOnly)
    \item \textbf{Scritture DB}: Nessuna diretta (gestite dal backend)
\end{itemize}

\textbf{Flusso di Autenticazione}:
\begin{enumerate}
    \item Utente inserisce credenziali
    \item Frontend chiama \texttt{authService.login()}
    \item Backend valida, genera JWT, setta cookie HttpOnly
    \item Frontend riceve risposta, aggiorna \texttt{AuthContext}
    \item Redirect alla dashboard
\end{enumerate}

\subsubsection{2. Dashboard (\texttt{features/dashboard})}

\textbf{Problema risolto}: Visualizzazione immediata delle metriche chiave e accesso rapido alle funzionalità principali.

\textbf{Componenti}:
\begin{itemize}
    \item \texttt{DashboardPage.tsx}: Pagina principale dashboard
    \item \texttt{QuickStats.tsx}: Statistiche rapide (ore totali, guadagni, obiettivi)
    \item \texttt{QuickActions.tsx}: Azioni rapide (nuovo progetto, nuovo worklog)
    \item \texttt{GoalsWidget.tsx}: Widget obiettivi in corso
    \item \texttt{AIInsightsWidget.tsx}: Insight generati da AI
    \item \texttt{WeeklyMiniChart.tsx}: Grafico settimanale compatto
\end{itemize}

\textbf{Servizi}:
\begin{itemize}
    \item \texttt{dashboardService.ts}: Chiamate API per dati dashboard
\end{itemize}

\textbf{Comunicazione}:
\begin{itemize}
    \item \textbf{API utilizzate}: \texttt{GET /api/dashboard/summary}, \texttt{GET /api/dashboard/quick-actions}
    \item \textbf{Dati scambiati}: Statistiche aggregate (ore, guadagni, progetti attivi, obiettivi)
    \item \textbf{Scritture DB}: Nessuna (solo letture)
\end{itemize}

\subsubsection{3. Obiettivi (\texttt{features/goals})}

\textbf{Problema risolto}: Gestione completa degli obiettivi personali con milestone, check-in, gamification e AI.

\textbf{Componenti}:
\begin{itemize}
    \item \texttt{GoalsPage.tsx}: Lista obiettivi con filtri
    \item \texttt{GoalDetailPage.tsx}: Dettaglio obiettivo (roadmap, journal, stats)
    \item \texttt{GoalWizard.tsx}: Wizard creazione obiettivo manuale
    \item \texttt{GoalWizardAI.tsx}: Wizard AI per creazione obiettivi da intent
    \item \texttt{ActivityHeatmap.tsx}: Heatmap attività (stile GitHub)
    \item \texttt{BurndownChart.tsx}: Grafico burndown per obiettivi target
    \item \texttt{MoodStats.tsx}: Statistiche mood dai check-in
\end{itemize}

\textbf{Context}:
\begin{itemize}
    \item \texttt{GoalsContext.tsx}: Stato globale obiettivi, CRUD operations
\end{itemize}

\textbf{Servizi}:
\begin{itemize}
    \item \texttt{goalsService.ts}: Chiamate API obiettivi
\end{itemize}

\textbf{Comunicazione}:
\begin{itemize}
    \item \textbf{API utilizzate}: 
    \begin{itemize}
        \item \texttt{GET /api/goals/goals}: Lista obiettivi
        \item \texttt{GET /api/goals/goals/:id}: Dettaglio obiettivo
        \item \texttt{POST /api/goals/goals}: Crea obiettivo
        \item \texttt{PUT /api/goals/goals/:id}: Aggiorna obiettivo
        \item \texttt{DELETE /api/goals/goals/:id}: Elimina obiettivo
        \item \texttt{POST /api/goals/goals/:goalId/quick-checkin}: Check-in rapido
        \item \texttt{POST /api/goals/goals/suggest}: Suggerimenti AI obiettivi
    \end{itemize}
    \item \textbf{Dati scambiati}: 
    \begin{itemize}
        \item \textbf{Input}: Titolo, categoria, tipo (TARGET/HABIT), milestone, action steps
        \item \textbf{Output}: Obiettivo completo con progresso, milestone completate, check-in history
    \end{itemize}
    \item \textbf{Scritture DB}: 
    \begin{itemize}
        \item \texttt{goals} collection: Inserimento/aggiornamento obiettivi
        \item \texttt{checkins} collection: Inserimento check-in con mood, note, progresso
    \end{itemize}
\end{itemize}

\subsubsection{4. Workspace (\texttt{features/workspace})}

\textbf{Problema risolto}: Gestione progetti lavorativi con work journal, todo list e tracking ore.

\textbf{Componenti}:
\begin{itemize}
    \item \texttt{WorkspacePage.tsx}: Pagina principale workspace
    \item \texttt{ProjectDetailPage.tsx}: Dettaglio progetto con entries e todos
    \item \texttt{ProjectCard.tsx}: Card progetto con statistiche
    \item \texttt{WorkEntryForm.tsx}: Form inserimento work entry
    \item \texttt{WorkEntryList.tsx}: Lista work entries con filtri
    \item \texttt{TodoList.tsx}: Lista TODO con drag \& drop
    \item \texttt{TimelineView.tsx}: Vista timeline cronologica
\end{itemize}

\textbf{Context}:
\begin{itemize}
    \item \texttt{WorkspaceContext.tsx}: Stato globale workspace (progetti, entries, todos)
\end{itemize}

\textbf{Servizi}:
\begin{itemize}
    \item \texttt{workspaceService.ts}: Chiamate API workspace
\end{itemize}

\textbf{Comunicazione}:
\begin{itemize}
    \item \textbf{API utilizzate}: 
    \begin{itemize}
        \item \texttt{GET /api/workspace/projects}: Lista progetti
        \item \texttt{POST /api/workspace/projects}: Crea progetto
        \item \texttt{GET /api/workspace/entries}: Lista work entries
        \item \texttt{POST /api/workspace/entries}: Crea work entry
        \item \texttt{GET /api/workspace/todos}: Lista TODO
        \item \texttt{POST /api/workspace/todos}: Crea TODO
    \end{itemize}
    \item \textbf{Dati scambiati}: 
    \begin{itemize}
        \item \textbf{Progetti}: Nome, descrizione, colore, stato
        \item \textbf{Work Entries}: Data, ora inizio/fine, descrizione, progetto
        \item \textbf{TODOs}: Titolo, descrizione, scadenza, priorità, stato
    \end{itemize}
    \item \textbf{Scritture DB}: 
    \begin{itemize}
        \item \texttt{workprojects} collection: Progetti workspace
        \item \texttt{workentries} collection: Work entries
        \item \texttt{worktodos} collection: TODO items
    \end{itemize}
\end{itemize}

\subsubsection{5. Tracker (\texttt{features/tracker})}

\textbf{Problema risolto}: Tracciamento ore lavorate (straordinari) con calcolo automatico guadagni.

\textbf{Componenti}:
\begin{itemize}
    \item \texttt{TimelinePage.tsx}: Timeline cronologica worklog
    \item \texttt{WorkLogForm.tsx}: Form inserimento worklog
    \item \texttt{WorkLogFormSmart.tsx}: Form intelligente con suggerimenti
    \item \texttt{WorkLogList.tsx}: Lista worklog con filtri mensili
    \item \texttt{ProjectSummary.tsx}: Riepilogo progetti con ore e guadagni
\end{itemize}

\textbf{Context}:
\begin{itemize}
    \item \texttt{WorkLogContext.tsx}: Stato globale worklog
\end{itemize}

\textbf{Comunicazione}:
\begin{itemize}
    \item \textbf{API utilizzate}: 
    \begin{itemize}
        \item \texttt{GET /api/worklogs}: Lista worklog
        \item \texttt{POST /api/worklogs}: Crea worklog
        \item \texttt{GET /api/worklogs/by-month/:year/:month}: Worklog per mese
        \item \texttt{GET /api/worklogs/totals}: Totali per periodo
    \end{itemize}
    \item \textbf{Dati scambiati}: 
    \begin{itemize}
        \item \textbf{Input}: projectId, data, startTime, endTime, description
        \item \textbf{Output}: Worklog con durata calcolata, guadagno (rate × ore)
    \end{itemize}
    \item \textbf{Scritture DB}: 
    \begin{itemize}
        \item \texttt{worklogs} collection: Inserimento worklog con calcolo \texttt{durationMinutes}
    \end{itemize}
\end{itemize}

\subsubsection{6. Studio (\texttt{features/study})}

\textbf{Problema risolto}: Sistema di studio con flashcard e spaced repetition (algoritmo SM-2).

\textbf{Componenti}:
\begin{itemize}
    \item \texttt{StudyPage.tsx}: Pagina principale studio
    \item \texttt{SplitStudyPage.tsx}: Vista split screen (card + note)
    \item \texttt{DeckList.tsx}: Lista deck flashcard
    \item \texttt{CardView.tsx}: Vista singola card
    \item \texttt{SessionSummary.tsx}: Riepilogo sessione studio
\end{itemize}

\textbf{Comunicazione}:
\begin{itemize}
    \item \textbf{API utilizzate}: 
    \begin{itemize}
        \item \texttt{GET /api/study/decks}: Lista deck
        \item \texttt{POST /api/study/decks}: Crea deck
        \item \texttt{POST /api/study/decks/:id/review}: Submit review card
        \item \texttt{POST /api/study/decks/:id/session-complete}: Completa sessione
    \end{itemize}
    \item \textbf{Dati scambiati}: 
    \begin{itemize}
        \item \textbf{Deck}: Nome, descrizione, cards
        \item \textbf{Card}: Front, back, ease factor, interval, next review date
        \item \textbf{Review}: Quality (0-5), timestamp
    \end{itemize}
    \item \textbf{Scritture DB}: 
    \begin{itemize}
        \item \texttt{decks} collection: Deck e cards con spaced repetition data
    \end{itemize}
\end{itemize}

\subsubsection{7. Impostazioni (\texttt{features/settings})}

\textbf{Problema risolto}: Gestione profilo utente, preferenze, notifiche, export dati.

\textbf{Componenti}:
\begin{itemize}
    \item \texttt{SettingsPage.tsx}: Pagina principale impostazioni
    \item \texttt{ProfileSettings.tsx}: Modifica profilo
    \item \texttt{PreferencesSettings.tsx}: Preferenze UI (tema, lingua)
    \item \texttt{NotificationsSettings.tsx}: Gestione notifiche
    \item \texttt{ExportSettings.tsx}: Export dati (GDPR)
    \item \texttt{DeleteAccountSettings.tsx}: Eliminazione account
\end{itemize}

\textbf{Comunicazione}:
\begin{itemize}
    \item \textbf{API utilizzate}: 
    \begin{itemize}
        \item \texttt{GET /api/settings/profile}: Profilo utente
        \item \texttt{PUT /api/settings/profile}: Aggiorna profilo
        \item \texttt{GET /api/settings/preferences}: Preferenze
        \item \texttt{PUT /api/settings/preferences}: Aggiorna preferenze
        \item \texttt{GET /api/settings/export}: Export dati
        \item \texttt{DELETE /api/settings/account}: Elimina account
    \end{itemize}
    \item \textbf{Dati scambiati}: 
    \begin{itemize}
        \item \textbf{Profilo}: firstName, lastName, displayName, phone, bio, avatar
        \item \textbf{Preferenze}: theme, language, timezone, dateFormat
    \end{itemize}
    \item \textbf{Scritture DB}: 
    \begin{itemize}
        \item \texttt{users} collection: Aggiornamento \texttt{profile} e \texttt{preferences}
    \end{itemize}
\end{itemize}

\subsubsection{8. Analytics (\texttt{features/analytics})}

\textbf{Problema risolto}: Visualizzazione analisi produttività e insights.

\textbf{Componenti}:
\begin{itemize}
    \item \texttt{ProductivityChart.tsx}: Grafico produttività settimanale
    \item \texttt{ActivityFeed.tsx}: Feed attività recenti
\end{itemize}

\textbf{Comunicazione}:
\begin{itemize}
    \item \textbf{API utilizzate}: 
    \begin{itemize}
        \item \texttt{GET /api/analytics/weekly}: Dati settimanali
        \item \texttt{GET /api/analytics/insights}: Insight AI
    \end{itemize}
    \item \textbf{Scritture DB}: Nessuna (solo letture aggregate)
\end{itemize}

\subsection{Componenti Condivisi (\texttt{shared})}

\subsubsection{API Client (\texttt{shared/services/apiClient.ts})}

\textbf{Problema risolto}: Gestione centralizzata delle chiamate HTTP con refresh token automatico, gestione errori e prevenzione "effetto valanga".

\textbf{Caratteristiche}:
\begin{itemize}
    \item \textbf{Mutex Pattern}: Solo una richiesta di refresh token alla volta
    \item \textbf{Request Queueing}: Richieste in attesa durante refresh
    \item \textbf{Auto Refresh}: Refresh automatico su 401
    \item \textbf{Error Handling}: Toast automatici con messaggi user-friendly
    \item \textbf{Anti-Spam}: Prevenzione toast duplicati
\end{itemize}

\textbf{Flusso Refresh Token}:
\begin{enumerate}
    \item Richiesta API fallisce con 401
    \item Se refresh già in corso: metti richiesta in coda
    \item Se primo refresh: acquisisci lock, chiama \texttt{POST /api/auth/refresh}
    \item Se successo: sblocca coda, riprova tutte le richieste
    \item Se fallimento: rifiuta coda, mostra toast, logout
\end{enumerate}

\subsubsection{Layouts (\texttt{shared/layouts})}

\textbf{Componenti}:
\begin{itemize}
    \item \texttt{AppLayout.tsx}: Layout principale con sidebar, header, navigazione
    \item \texttt{AuthLayout.tsx}: Layout autenticazione (centrato, senza sidebar)
\end{itemize}

\subsubsection{Toast System (\texttt{shared/components/toast})}

\textbf{Problema risolto}: Sistema di notifiche centralizzato con deduplicazione.

\textbf{Componenti}:
\begin{itemize}
    \item \texttt{ToastContext.tsx}: Context provider per toast
    \item \texttt{toastEvents.ts}: Event emitter per toast da codice non-React
    \item \texttt{types.ts}: TypeScript types per toast
\end{itemize}

\section{Componenti Backend}

Il backend segue un'architettura \textbf{MVC (Model-View-Controller)} con separazione delle responsabilità.

\subsection{Struttura Backend}

\begin{itemize}
    \item \texttt{models/}: Schemi Mongoose (User, Goal, WorkLog, Project, etc.)
    \item \texttt{controllers/}: Logica di gestione richieste HTTP
    \item \texttt{services/}: Business logic (separata dai controller)
    \item \texttt{routes/}: Definizione endpoint API
    \item \texttt{middleware/}: Autenticazione, error handling, rate limiting, tenant context
    \item \texttt{config/}: Configurazione sicurezza, Redis, CORS
    \item \texttt{utils/}: Utility (AppError, validators)
\end{itemize}

\subsection{Modelli Database (MongoDB)}

\subsubsection{User Model}

\textbf{Collection}: \texttt{users}

\textbf{Campi Principali}:
\begin{itemize}
    \item \texttt{email}: String (unique, indexed)
    \item \texttt{password}: String (hashed con Argon2id, select: false)
    \item \texttt{profile}: Object (firstName, lastName, displayName, phone, bio, avatar, company, jobTitle, location, website)
    \item \texttt{preferences}: Object (theme, language, timezone, dateFormat, notifications)
    \item \texttt{refreshTokens}: Array (sessioni multi-device)
    \item \texttt{gracePeriodTokens}: Array (token temporanei per race conditions)
    \item \texttt{emailVerified}: Boolean
    \item \texttt{emailVerificationToken}: String
    \item \texttt{passwordResetToken}: String
    \item \texttt{passwordResetExpires}: Date
    \item \texttt{isActive}: Boolean
    \item \texttt{lastLoginAt}: Date
    \item \texttt{loginAttempts}: Number
    \item \texttt{lockUntil}: Date
\end{itemize}

\textbf{Metodi}:
\begin{itemize}
    \item \texttt{comparePassword()}: Verifica password
    \item \texttt{addSession()}: Aggiunge sessione refresh token
    \item \texttt{removeSessionByHash()}: Rimuove sessione specifica
    \item \texttt{removeAllSessions()}: Rimuove tutte le sessioni
    \item \texttt{addToGracePeriod()}: Aggiunge token a grace period
    \item \texttt{removeFromGracePeriod()}: Rimuove token scaduti
\end{itemize}

\textbf{Scritture DB}:
\begin{itemize}
    \item \textbf{Registrazione}: Inserimento nuovo documento User
    \item \textbf{Login}: Aggiornamento \texttt{refreshTokens}, \texttt{lastLoginAt}
    \item \textbf{Refresh Token}: Aggiornamento \texttt{refreshTokens} (rotazione), \texttt{gracePeriodTokens}
    \item \textbf{Logout}: Aggiornamento \texttt{refreshTokens} (rimozione sessione)
    \item \textbf{Cambio Password}: Aggiornamento \texttt{password}, reset \texttt{refreshTokens}
\end{itemize}

\subsubsection{Goal Model}

\textbf{Collection}: \texttt{goals}

\textbf{Campi Principali}:
\begin{itemize}
    \item \texttt{user}: ObjectId (ref: User, multi-tenancy)
    \item \texttt{title}: String
    \item \texttt{category}: String (enum: HEALTH, CAREER, FINANCE, LEARNING, RELATIONSHIPS, PERSONAL, CREATIVE, TRAVEL, OTHER)
    \item \texttt{type}: String (enum: TARGET, HABIT)
    \item \texttt{targetValue}: Number (per TARGET)
    \item \texttt{currentValue}: Number (per TARGET)
    \item \texttt{unit}: String (per TARGET: EUR, USD, KG, KM, etc.)
    \item \texttt{frequency}: Object (per HABIT: timesPerWeek, timesPerMonth)
    \item \texttt{milestones}: Array (titolo, note, isCompleted, weight, completedAt, actionSteps)
    \item \texttt{deadline}: Date
    \item \texttt{status}: String (enum: ACTIVE, COMPLETED, PAUSED, CANCELLED)
    \item \texttt{progress}: Number (0-100)
    \item \texttt{aiGenerated}: Boolean
    \item \texttt{aiPlan}: Object (per obiettivi generati da AI)
\end{itemize}

\textbf{Scritture DB}:
\begin{itemize}
    \item \textbf{Creazione}: Inserimento nuovo documento Goal
    \item \textbf{Aggiornamento}: Update milestone, progress, status
    \item \textbf{Check-in}: Inserimento in \texttt{checkins} collection (collegato via \texttt{goalId})
\end{itemize}

\subsubsection{WorkLog Model}

\textbf{Collection}: \texttt{worklogs}

\textbf{Campi Principali}:
\begin{itemize}
    \item \texttt{user}: ObjectId (ref: User, multi-tenancy)
    \item \texttt{projectId}: ObjectId (ref: Project)
    \item \texttt{date}: String (YYYY-MM-DD)
    \item \texttt{startTime}: String (HH:mm)
    \item \texttt{endTime}: String (HH:mm)
    \item \texttt{description}: String
    \item \texttt{durationMinutes}: Number (calcolato automaticamente)
\end{itemize}

\textbf{Indici}:
\begin{itemize}
    \item \texttt{\{user: 1, date: -1\}}: Query per mese
    \item \texttt{\{user: 1, projectId: 1, date: -1\}}: Query per progetto
\end{itemize}

\textbf{Scritture DB}:
\begin{itemize}
    \item \textbf{Creazione}: Inserimento worklog con calcolo \texttt{durationMinutes}
    \item \textbf{Aggiornamento}: Modifica date/times, ricalcolo durata
    \item \textbf{Eliminazione}: Rimozione worklog
\end{itemize}

\subsubsection{Project Model}

\textbf{Collection}: \texttt{projects}

\textbf{Campi Principali}:
\begin{itemize}
    \item \texttt{user}: ObjectId (ref: User, multi-tenancy)
    \item \texttt{name}: String
    \item \texttt{code}: String (codice commessa)
    \item \texttt{description}: String
    \item \texttt{rate}: Number (tariffa oraria)
    \item \texttt{estimatedHours}: Number
    \item \texttt{progress}: Number (0-100)
    \item \texttt{status}: String (enum: ACTIVE, COMPLETED, ON_HOLD, CANCELLED)
    \item \texttt{color}: String (hex color)
\end{itemize}

\textbf{Scritture DB}:
\begin{itemize}
    \item \textbf{Creazione}: Inserimento progetto
    \item \textbf{Aggiornamento}: Modifica rate, status, progress
    \item \textbf{Eliminazione}: Solo se non ha worklog associati
\end{itemize}

\subsubsection{Altri Modelli}

\begin{itemize}
    \item \texttt{CheckIn}: Check-in obiettivi (mood, note, progresso)
    \item \texttt{Deck}: Deck flashcard per studio
    \item \texttt{WorkEntry}: Work entry nel workspace
    \item \texttt{WorkTodo}: TODO items nel workspace
    \item \texttt{UserActivity}: Log attività utente (analytics)
\end{itemize}

\subsection{Controllers}

I controller gestiscono le richieste HTTP e delegano la business logic ai servizi.

\subsubsection{Auth Controller}

\textbf{File}: \texttt{controllers/authController.js}

\textbf{Metodi}:
\begin{itemize}
    \item \texttt{register()}: Registrazione nuovo utente
    \item \texttt{login()}: Login utente
    \item \texttt{logout()}: Logout (rimozione sessione)
    \item \texttt{refresh()}: Refresh access token
    \item \texttt{check()}: Verifica sessione corrente
    \item \texttt{verifyEmail()}: Verifica email con token
    \item \texttt{forgotPassword()}: Richiesta reset password
    \item \texttt{resetPassword()}: Reset password con token
    \item \texttt{changePassword()}: Cambio password (autenticato)
\end{itemize}

\textbf{Flusso Login}:
\begin{enumerate}
    \item Riceve email/password
    \item Valida con Zod
    \item Chiama \texttt{authService.login()}
    \item Genera JWT access token (15 min)
    \item Genera refresh token (7 giorni)
    \item Hash refresh token, salva in \texttt{refreshTokens} array
    \item Setta cookie HttpOnly con tokens
    \item Ritorna user (senza password)
\end{enumerate}

\subsubsection{Goal Controller}

\textbf{File}: \texttt{controllers/goalController.js}

\textbf{Metodi}:
\begin{itemize}
    \item \texttt{getGoals()}: Lista obiettivi utente
    \item \texttt{getGoalById()}: Dettaglio obiettivo
    \item \texttt{createGoal()}: Crea obiettivo
    \item \texttt{updateGoal()}: Aggiorna obiettivo
    \item \texttt{deleteGoal()}: Elimina obiettivo
    \item \texttt{toggleMilestone()}: Completa/riapre milestone
    \item \texttt{quickCheckIn()}: Check-in rapido
    \item \texttt{getCheckIns()}: Lista check-in obiettivo
    \item \texttt{getStats()}: Statistiche obiettivi
\end{itemize}

\subsubsection{Altri Controllers}

\begin{itemize}
    \item \texttt{workspaceController.js}: Gestione workspace (progetti, entries, todos)
    \item \texttt{dashboardController.js}: Dati dashboard
    \item \texttt{analyticsController.js}: Analytics e insights
    \item \texttt{settingsController.js}: Impostazioni utente
    \item \texttt{studyController.js}: Studio con flashcard
\end{itemize}

\subsection{Services}

I servizi contengono la business logic pura, separata dalla gestione HTTP.

\subsubsection{Auth Service}

\textbf{File}: \texttt{services/authService.js}

\textbf{Metodi}:
\begin{itemize}
    \item \texttt{register()}: Crea utente, hash password, genera email verification token
    \item \texttt{login()}: Verifica credenziali, genera tokens, gestisce sessioni multi-device
    \item \texttt{refreshAccessToken()}: Ruota refresh token, genera nuovo access token
    \item \texttt{logout()}: Rimuove sessione specifica o tutte le sessioni
    \item \texttt{changePassword()}: Cambia password, invalida tutte le sessioni
    \item \texttt{addToBlacklist()}: Aggiunge token a blacklist Redis
    \item \texttt{isTokenBlacklisted()}: Verifica se token è in blacklist
\end{itemize}

\textbf{Scritture DB}:
\begin{itemize}
    \item \textbf{Register}: \texttt{User.create()} con password hashata
    \item \textbf{Login}: \texttt{User.updateOne()} con \texttt{\$push} per \texttt{refreshTokens} (FIFO limit 10)
    \item \textbf{Refresh}: \texttt{User.findOneAndUpdate()} con aggregation pipeline (rimozione vecchio token, inserimento nuovo, FIFO)
    \item \textbf{Logout}: \texttt{User.updateOne()} con \texttt{\$pull} per rimuovere sessione
\end{itemize}

\subsubsection{Goal Service}

\textbf{File}: \texttt{services/goalService.js}

\textbf{Metodi}:
\begin{itemize}
    \item \texttt{create()}: Crea obiettivo con validazione
    \item \texttt{update()}: Aggiorna obiettivo
    \item \texttt{calculateProgress()}: Calcola progresso automatico
    \item \texttt{getStats()}: Statistiche aggregate
\end{itemize}

\textbf{Scritture DB}:
\begin{itemize}
    \item \textbf{Creazione}: \texttt{Goal.create()} con \texttt{user} dal \texttt{tenantScope}
    \item \textbf{Aggiornamento}: \texttt{Goal.updateOne()} con filtro \texttt{user}
    \item \textbf{Check-in}: \texttt{CheckIn.create()} collegato a goal
\end{itemize}

\subsubsection{WorkLog Service}

\textbf{File}: \texttt{services/workLogService.js}

\textbf{Metodi}:
\begin{itemize}
    \item \texttt{create()}: Crea worklog, verifica \texttt{projectId} appartiene a utente, calcola \texttt{durationMinutes}
    \item \texttt{findByMonth()}: Query worklog per mese
    \item \texttt{getTotalsByPeriod()}: Aggregazione totali (ore, guadagni) per periodo
\end{itemize}

\textbf{Scritture DB}:
\begin{itemize}
    \item \textbf{Creazione}: \texttt{WorkLog.create()} con calcolo durata
    \item \textbf{Query}: \texttt{WorkLog.aggregate()} per statistiche
\end{itemize}

\subsubsection{Base Service}

\textbf{File}: \texttt{services/BaseService.js}

\textbf{Pattern}: Classe base per tutti i servizi con metodi comuni:
\begin{itemize}
    \item \texttt{find()}: Lista con filtro \texttt{tenantScope}
    \item \texttt{findById()}: Dettaglio con verifica ownership
    \item \texttt{create()}: Creazione con \texttt{user} automatico
    \item \texttt{update()}: Update con verifica ownership
    \item \texttt{delete()}: Delete con verifica ownership
\end{itemize}

\textbf{Multi-Tenancy}: Tutti i metodi filtrano automaticamente per \texttt{req.tenantScope.userId}.

\subsection{Routes}

Le routes definiscono gli endpoint API e collegano middleware, controller e servizi.

\subsubsection{Auth Routes}

\textbf{File}: \texttt{routes/auth.js}

\textbf{Endpoints}:
\begin{itemize}
    \item \texttt{POST /api/auth/register}: Registrazione
    \item \texttt{POST /api/auth/login}: Login
    \item \texttt{POST /api/auth/logout}: Logout
    \item \texttt{POST /api/auth/refresh}: Refresh token
    \item \texttt{GET /api/auth/check}: Verifica sessione
    \item \texttt{GET /api/auth/verify-email/:token}: Verifica email
    \item \texttt{POST /api/auth/forgot-password}: Richiesta reset password
    \item \texttt{POST /api/auth/reset-password}: Reset password
    \item \texttt{PUT /api/auth/change-password}: Cambio password
\end{itemize}

\subsubsection{Goals Routes}

\textbf{File}: \texttt{routes/goals.js}

\textbf{Endpoints}:
\begin{itemize}
    \item \texttt{GET /api/goals/goals}: Lista obiettivi
    \item \texttt{GET /api/goals/goals/:id}: Dettaglio obiettivo
    \item \texttt{POST /api/goals/goals}: Crea obiettivo
    \item \texttt{PUT /api/goals/goals/:id}: Aggiorna obiettivo
    \item \texttt{DELETE /api/goals/goals/:id}: Elimina obiettivo
    \item \texttt{PATCH /api/goals/goals/:id/milestones/:milestoneId/toggle}: Toggle milestone
    \item \texttt{POST /api/goals/goals/:goalId/quick-checkin}: Check-in rapido
    \item \texttt{GET /api/goals/goals/:goalId/checkins}: Lista check-in
    \item \texttt{POST /api/goals/goals/suggest}: Suggerimenti AI
    \item \texttt{GET /api/goals/goals-stats}: Statistiche obiettivi
\end{itemize}

\subsubsection{API Routes (Workspace)}

\textbf{File}: \texttt{routes/api.js}

\textbf{Endpoints}:
\begin{itemize}
    \item \texttt{GET /api/projects}: Lista progetti
    \item \texttt{POST /api/projects}: Crea progetto
    \item \texttt{GET /api/worklogs}: Lista worklog
    \item \texttt{POST /api/worklogs}: Crea worklog
    \item \texttt{GET /api/workspace/projects}: Lista progetti workspace
    \item \texttt{POST /api/workspace/entries}: Crea work entry
    \item \texttt{GET /api/workspace/todos}: Lista TODO
    \item \texttt{POST /api/workspace/todos}: Crea TODO
\end{itemize}

\subsection{Middleware}

\subsubsection{Auth Middleware}

\textbf{File}: \texttt{middleware/auth.js}

\textbf{Funzione}: Verifica JWT access token, estrae payload, aggiunge \texttt{req.user}.

\textbf{Flusso}:
\begin{enumerate}
    \item Estrae token da cookie o header
    \item Verifica firma JWT
    \item Verifica scadenza
    \item Verifica blacklist Redis (se token revocato)
    \item Verifica utente attivo nel DB
    \item Aggiunge \texttt{req.user} alla request
\end{enumerate}

\subsubsection{Tenant Context Middleware}

\textbf{File}: \texttt{middleware/tenantContext.js}

\textbf{Funzione}: Inietta \texttt{req.tenantScope} con \texttt{userId} per multi-tenancy.

\textbf{Utilizzo}: Tutti i servizi usano \texttt{req.tenantScope} per filtrare automaticamente i dati.

\subsubsection{Rate Limiter Middleware}

\textbf{File}: \texttt{middleware/rateLimiter.js}

\textbf{Funzione}: Limita richieste per IP/utente per prevenire brute force e DDoS.

\textbf{Configurazione}:
\begin{itemize}
    \item \textbf{Store}: Redis (distribuito) o MemoryStore (locale)
    \item \textbf{Limiti}: 100 richieste/15 min (generale), 5 tentativi/15 min (login)
\end{itemize}

\subsubsection{Error Handler Middleware}

\textbf{File}: \texttt{middleware/errorHandler.js}

\textbf{Funzione}: Gestione centralizzata errori, logging, risposte user-friendly.

\textbf{Caratteristiche}:
\begin{itemize}
    \item Distingue errori operazionali (4xx) da bug (5xx)
    \item Log dettagliato in sviluppo, messaggi user-friendly in produzione
    \item Silenzia 401 su \texttt{/auth/check} (utente anonimo normale)
    \item Trasforma errori Mongoose/JWT in messaggi leggibili
\end{itemize}

\section{Comunicazione tra Componenti}

\subsection{Flusso Richiesta API}

\begin{enumerate}
    \item \textbf{Frontend}: Componente chiama \texttt{apiClient.get('/api/goals/goals')}
    \item \textbf{API Client}: Interceptor aggiunge cookie, gestisce refresh token
    \item \textbf{Network}: Richiesta HTTP a backend
    \item \textbf{Backend}: Express riceve richiesta
    \item \textbf{Middleware Chain}:
    \begin{itemize}
        \item Helmet (security headers)
        \item CORS
        \item Rate Limiter
        \item Body Parser
        \item Auth Middleware (verifica JWT)
        \item Tenant Context (inietta userId)
    \end{itemize}
    \item \textbf{Route}: \texttt{GET /api/goals/goals} → \texttt{goalsController.getGoals()}
    \item \textbf{Controller}: Chiama \texttt{goalService.find(req.tenantScope)}
    \item \textbf{Service}: Query MongoDB con filtro \texttt{user: req.tenantScope.userId}
    \item \textbf{Database}: Ritorna documenti Goal
    \item \textbf{Service}: Ritorna array obiettivi
    \item \textbf{Controller}: \texttt{res.json(\{success: true, data: goals\})}
    \item \textbf{Frontend}: Riceve risposta, aggiorna Context, re-render UI
\end{enumerate}

\subsection{Flusso Autenticazione}

\begin{enumerate}
    \item \textbf{Login}: Frontend invia email/password
    \item \textbf{Backend}: Valida, hash password, genera JWT
    \item \textbf{Backend}: Salva refresh token hash in \texttt{refreshTokens} array
    \item \textbf{Backend}: Setta cookie HttpOnly con tokens
    \item \textbf{Frontend}: Riceve risposta, aggiorna \texttt{AuthContext}
    \item \textbf{Richieste successive}: Cookie inviati automaticamente
    \item \textbf{401 su richiesta}: Frontend interceptor cattura, fa refresh automatico
    \item \textbf{Refresh}: Backend ruota refresh token, genera nuovo access token
    \item \textbf{Frontend}: Riprova richiesta originale con nuovo token
\end{enumerate}

\subsection{Flusso Multi-Device}

\begin{enumerate}
    \item \textbf{Login Device A}: Backend aggiunge sessione a \texttt{refreshTokens} array
    \item \textbf{Login Device B}: Backend aggiunge altra sessione (FIFO: max 10)
    \item \textbf{Refresh Device A}: Backend trova sessione, ruota token, aggiorna array
    \item \textbf{Logout Device B}: Backend rimuove solo sessione Device B
    \item \textbf{Device A}: Continua a funzionare (sessione indipendente)
\end{enumerate}

\section{Scritture Database}

\subsection{Operazioni CRUD Principali}

\subsubsection{User}

\begin{itemize}
    \item \textbf{CREATE}: \texttt{User.create(\{email, password: hashed\})}
    \item \textbf{UPDATE}: \texttt{User.updateOne(\{_id\}, \{\$push: \{refreshTokens\}\})}
    \item \textbf{UPDATE}: \texttt{User.updateOne(\{_id\}, \{\$pull: \{refreshTokens: \{hash\}\}\})}
    \item \textbf{UPDATE}: \texttt{User.updateOne(\{_id\}, \{\$set: \{profile, preferences\}\})}
\end{itemize}

\subsubsection{Goal}

\begin{itemize}
    \item \textbf{CREATE}: \texttt{Goal.create(\{user, title, category, type, milestones\})}
    \item \textbf{UPDATE}: \texttt{Goal.updateOne(\{_id, user\}, \{\$set: \{progress, status\}\})}
    \item \textbf{UPDATE}: \texttt{Goal.updateOne(\{_id, user\}, \{\$set: \{"milestones.\$[elem].isCompleted": true\}\})}
    \item \textbf{DELETE}: \texttt{Goal.deleteOne(\{_id, user\})}
\end{itemize}

\subsubsection{WorkLog}

\begin{itemize}
    \item \textbf{CREATE}: \texttt{WorkLog.create(\{user, projectId, date, startTime, endTime, durationMinutes\})}
    \item \textbf{UPDATE}: \texttt{WorkLog.updateOne(\{_id, user\}, \{\$set: \{date, startTime, endTime, durationMinutes\}\})}
    \item \textbf{DELETE}: \texttt{WorkLog.deleteOne(\{_id, user\})}
\end{itemize}

\subsection{Operazioni Atomiche}

\subsubsection{Refresh Token Rotation}

\textbf{Problema}: Evitare race conditions su \texttt{refreshTokens} array.

\textbf{Soluzione}: Aggregation pipeline atomica:

\begin{lstlisting}
User.findOneAndUpdate(
    { _id: userId, 'refreshTokens.hash': oldHash },
    [
        { $set: { refreshTokens: { $filter: { input: '$refreshTokens', cond: { $ne: ['$$this.hash', oldHash] } } } } },
        { $set: { refreshTokens: { $concatArrays: ['$refreshTokens', [newSession]] } } },
        { $set: { refreshTokens: { $slice: [{ $sortArray: { input: '$refreshTokens', sortBy: { lastUsedAt: -1 } } }, MAX_SESSIONS] } } }
    ]
)
\end{lstlisting}

\subsubsection{FIFO Session Limit}

\textbf{Problema}: Array \texttt{refreshTokens} può crescere infinitamente.

\textbf{Soluzione}: Limite FIFO (First In, First Out) con aggregation pipeline:

\begin{lstlisting}
User.updateOne(
    { _id: userId },
    [
        { $set: { refreshTokens: { $concatArrays: ['$refreshTokens', [newSession]] } } },
        { $set: { refreshTokens: { $slice: [{ $sortArray: { input: '$refreshTokens', sortBy: { lastUsedAt: 1 } } }, -MAX_SESSIONS] } } }
    ]
)
\end{lstlisting}

\section{API Endpoints Completi}

\subsection{Autenticazione}

\begin{longtable}{|p{3cm}|p{2cm}|p{4cm}|p{6cm}|}
\hline
\textbf{Endpoint} & \textbf{Method} & \textbf{Input} & \textbf{Output} \\
\hline
\endhead
\hline
\endfoot
\hline
\endlastfoot
\texttt{/api/auth/register} & POST & \texttt{\{email, password\}} & \texttt{\{success, data: \{user\}\}} \\
\texttt{/api/auth/login} & POST & \texttt{\{email, password\}} & \texttt{\{success, data: \{user\}\}} (cookie: accessToken, refreshToken) \\
\texttt{/api/auth/logout} & POST & Cookie: refreshToken & \texttt{\{success, message\}} \\
\texttt{/api/auth/refresh} & POST & Cookie: refreshToken & \texttt{\{success\}} (cookie: accessToken, refreshToken) \\
\texttt{/api/auth/check} & GET & Cookie: accessToken & \texttt{\{success, data: \{user, isAuthenticated\}\}} \\
\texttt{/api/auth/verify-email/:token} & GET & Token in URL & \texttt{\{success, message\}} \\
\texttt{/api/auth/forgot-password} & POST & \texttt{\{email\}} & \texttt{\{success, message\}} \\
\texttt{/api/auth/reset-password} & POST & \texttt{\{token, password\}} & \texttt{\{success, message\}} \\
\texttt{/api/auth/change-password} & PUT & \texttt{\{currentPassword, newPassword\}} & \texttt{\{success, message\}} \\
\hline
\end{longtable}

\subsection{Obiettivi}

\begin{longtable}{|p{3cm}|p{2cm}|p{4cm}|p{6cm}|}
\hline
\textbf{Endpoint} & \textbf{Method} & \textbf{Input} & \textbf{Output} \\
\hline
\endhead
\hline
\endfoot
\hline
\endlastfoot
\texttt{/api/goals/goals} & GET & Query: \texttt{?status=ACTIVE} & \texttt{\{success, data: [goals]\}} \\
\texttt{/api/goals/goals/:id} & GET & - & \texttt{\{success, data: goal\}} \\
\texttt{/api/goals/goals} & POST & \texttt{\{title, category, type, milestones\}} & \texttt{\{success, data: goal\}} \\
\texttt{/api/goals/goals/:id} & PUT & \texttt{\{title, progress, status\}} & \texttt{\{success, data: goal\}} \\
\texttt{/api/goals/goals/:id} & DELETE & - & \texttt{\{success, message\}} \\
\texttt{/api/goals/goals/:id/milestones/:milestoneId/toggle} & PATCH & - & \texttt{\{success, data: goal\}} \\
\texttt{/api/goals/goals/:goalId/quick-checkin} & POST & \texttt{\{mood, notes, progress\}} & \texttt{\{success, data: checkin\}} \\
\texttt{/api/goals/goals/:goalId/checkins} & GET & - & \texttt{\{success, data: [checkins]\}} \\
\texttt{/api/goals/goals/suggest} & POST & \texttt{\{intent, category, intensity\}} & \texttt{\{success, data: \{plan\}\}} \\
\texttt{/api/goals/goals-stats} & GET & - & \texttt{\{success, data: \{stats\}\}} \\
\hline
\end{longtable}

\subsection{Workspace}

\begin{longtable}{|p{3cm}|p{2cm}|p{4cm}|p{6cm}|}
\hline
\textbf{Endpoint} & \textbf{Method} & \textbf{Input} & \textbf{Output} \\
\hline
\endhead
\hline
\endfoot
\hline
\endlastfoot
\texttt{/api/workspace/projects} & GET & - & \texttt{\{success, data: [projects]\}} \\
\texttt{/api/workspace/projects} & POST & \texttt{\{name, description, color\}} & \texttt{\{success, data: project\}} \\
\texttt{/api/workspace/entries} & GET & Query: \texttt{?projectId=xxx} & \texttt{\{success, data: [entries]\}} \\
\texttt{/api/workspace/entries} & POST & \texttt{\{projectId, date, startTime, endTime, description\}} & \texttt{\{success, data: entry\}} \\
\texttt{/api/workspace/todos} & GET & Query: \texttt{?projectId=xxx} & \texttt{\{success, data: [todos]\}} \\
\texttt{/api/workspace/todos} & POST & \texttt{\{projectId, title, description, dueDate, priority\}} & \texttt{\{success, data: todo\}} \\
\texttt{/api/workspace/todos/:id/complete} & PATCH & - & \texttt{\{success, data: todo\}} \\
\hline
\end{longtable}

\subsection{Dashboard e Analytics}

\begin{longtable}{|p{3cm}|p{2cm}|p{4cm}|p{6cm}|}
\hline
\textbf{Endpoint} & \textbf{Method} & \textbf{Input} & \textbf{Output} \\
\hline
\endhead
\hline
\endfoot
\hline
\endlastfoot
\texttt{/api/dashboard/summary} & GET & - & \texttt{\{success, data: \{hours, earnings, goals, projects\}\}} \\
\texttt{/api/dashboard/quick-actions} & GET & - & \texttt{\{success, data: [actions]\}} \\
\texttt{/api/analytics/weekly} & GET & - & \texttt{\{success, data: [dailyStats]\}} \\
\texttt{/api/analytics/insights} & GET & - & \texttt{\{success, data: [insights]\}} \\
\hline
\end{longtable}

\section{Pattern e Architetture}

\subsection{Multi-Tenancy}

\textbf{Implementazione}: Plugin Mongoose che aggiunge automaticamente filtro \texttt{user} a tutte le query.

\textbf{Vantaggi}:
\begin{itemize}
    \item Isolamento completo dati per utente
    \item Nessun rischio di accesso cross-user
    \item Query automaticamente filtrate
\end{itemize}

\textbf{Utilizzo}: Tutti i modelli (tranne User) hanno campo \texttt{user: ObjectId} e usano \texttt{multiTenancyPlugin}.

\subsection{Refresh Token Rotation}

\textbf{Pattern}: Ad ogni refresh, il vecchio token viene invalidato e ne viene generato uno nuovo.

\textbf{Vantaggi}:
\begin{itemize}
    \item Rilevamento furto token (se vecchio token viene riutilizzato)
    \item Limite danno in caso di compromissione
\end{itemize}

\textbf{Grace Period}: Vecchi token rimangono validi per 10-30 secondi per gestire race conditions.

\subsection{Mutex Pattern (Frontend)}

\textbf{Problema}: Multiple richieste simultanee che falliscono con 401 causano multiple chiamate refresh.

\textbf{Soluzione}: Solo una richiesta può fare refresh alla volta, le altre aspettano in coda.

\textbf{Implementazione}: Flag \texttt{isRefreshing} + array \texttt{failedQueue} in \texttt{apiClient.ts}.

\subsection{Error Handling Centralizzato}

\textbf{Backend}: Middleware \texttt{errorHandler} cattura tutti gli errori, logga dettagli, invia messaggi user-friendly.

\textbf{Frontend}: Interceptor Axios cattura errori, estrae messaggio backend, mostra toast automatico.

\subsection{Servizi Aggiuntivi}

\subsubsection{Activity Service}

\textbf{File}: \texttt{services/activityService.js}

\textbf{Problema risolto}: Tracciamento attività utente per gamification e analytics.

\textbf{Funzionalità}:
\begin{itemize}
    \item Calcolo XP (Experience Points) per azioni utente
    \item Sistema di livelli basato su XP
    \item Gestione streak (giorni consecutivi di attività)
    \item Log attività per analytics
\end{itemize}

\textbf{Scritture DB}:
\begin{itemize}
    \item \texttt{useractivities} collection: Log attività con timestamp, tipo, XP guadagnato
    \item \texttt{users} collection: Aggiornamento \texttt{xp}, \texttt{level}, \texttt{streak}
\end{itemize}

\subsubsection{Insight Service}

\textbf{File}: \texttt{services/insightService.js}

\textbf{Problema risolto}: Generazione insight AI per produttività e obiettivi.

\textbf{Funzionalità}:
\begin{itemize}
    \item Analisi pattern produttività
    \item Suggerimenti personalizzati
    \item Rilevamento anomalie (es. calo produttività)
\end{itemize}

\textbf{Comunicazione}: Utilizza OpenAI API per generare insight basati su dati utente aggregati.

\subsubsection{Email Service}

\textbf{File}: \texttt{services/emailService.js}

\textbf{Problema risolto}: Invio email per verifica account, reset password, notifiche.

\textbf{Funzionalità}:
\begin{itemize}
    \item Email verifica account
    \item Email reset password
    \item Notifiche obiettivi (opzionale)
\end{itemize}

\textbf{Configurazione}: Utilizza servizio email esterno (es. SendGrid, AWS SES) configurato via variabili ambiente.

\subsubsection{Vector Store Service}

\textbf{File}: \texttt{services/vectorStoreService.js}

\textbf{Problema risolto}: Storage e ricerca vettoriale per contenuti PDF e documenti (studio).

\textbf{Funzionalità}:
\begin{itemize}
    \item Embedding di testo con OpenAI
    \item Storage vettori in MongoDB
    \item Ricerca semantica per contenuti simili
\end{itemize}

\textbf{Utilizzo}: Sistema di studio con flashcard, estrazione contenuti da PDF.

\subsubsection{AI Goal Service}

\textbf{File}: \texttt{services/aiGoalService.js}

\textbf{Problema risolto}: Generazione automatica di piani obiettivi da intent utente.

\textbf{Funzionalità}:
\begin{itemize}
    \item Analisi intent utente (testo libero)
    \item Generazione milestone e action steps
    \item Calcolo intensità e difficoltà
\end{itemize}

\textbf{Comunicazione}: Utilizza OpenAI API per generare piani strutturati da intent naturali.

\subsubsection{Spaced Repetition Algorithms}

\textbf{File}: \texttt{services/spacedRepetitionAlgorithms.js}

\textbf{Problema risolto}: Algoritmi di spaced repetition per ottimizzare apprendimento flashcard.

\textbf{Algoritmi}:
\begin{itemize}
    \item \textbf{SM-2}: Algoritmo classico SuperMemo 2
    \item \textbf{FSRS}: Free Spaced Repetition Scheduler (moderno)
\end{itemize}

\textbf{Calcoli}:
\begin{itemize}
    \item Ease factor (facilità di richiamo)
    \item Interval (giorni fino prossima review)
    \item Next review date
\end{itemize}

\section{Conclusioni}

Il sistema ExtraTracker è progettato con un'architettura modulare, scalabile e sicura. La separazione tra frontend e backend, l'uso di pattern consolidati (MVC, Multi-Tenancy, Mutex) e l'implementazione di sicurezza a più livelli (JWT, rate limiting, CORS, Helmet) garantiscono un'applicazione robusta e pronta per la produzione.

\subsection{Riepilogo Componenti}

\begin{itemize}
    \item \textbf{Frontend}: 8 features principali + componenti condivisi
    \item \textbf{Backend}: 8 controllers + 15+ services + 10+ models
    \item \textbf{API Endpoints}: 50+ endpoint REST
    \item \textbf{Database}: 10+ collections MongoDB
    \item \textbf{Sicurezza}: JWT, Argon2id, Rate Limiting, CORS, Helmet, Multi-Tenancy
    \item \textbf{Scalabilità}: Redis per cache, pattern Mutex, operazioni atomiche MongoDB
\end{itemize}

\subsection{Pattern Architetturali Utilizzati}

\begin{itemize}
    \item \textbf{MVC (Model-View-Controller)}: Separazione responsabilità backend
    \item \textbf{Repository Pattern}: BaseService astrae accesso dati
    \item \textbf{Template Method Pattern}: BaseService definisce scheletro algoritmi
    \item \textbf{Multi-Tenancy}: Isolamento dati per utente
    \item \textbf{Mutex Pattern}: Prevenzione race conditions frontend
    \item \textbf{Error Handling Centralizzato}: Gestione errori unificata
    \item \textbf{Service Layer}: Business logic separata da HTTP
\end{itemize}

\vspace{1cm}

\textit{Documento generato il \today}

\end{document}
